\subsection{Equicontinuity}

% **Role:** convergence of functions and analytic approximation.
%
% Equicontinuity, Arzela--Ascoli, compactness in function spaces, Banach-space preview.


Function-space compactness treats functions as points. Supremum metrics preview Banach spaces, Baire category describes generic behavior, and equicontinuity leads to Arzela--Ascoli.

\begin{theorem} \label{thm:complete-supremum-metric-space} \hlt{Completeness of the Supremum Metric}\\
The metric space $C_b(X)$ with metric $d(f,g)=\norm{f-g}_{\infty}$ is complete.
\begin{proof}
Let $(f_n)$ be Cauchy in $C_b(X)$. Then Theorem \ref{thm:uniform-cauchy-criterion} gives a function $f$ with $f_n\rightrightarrows f$. By Theorem \ref{thm:uniform-convergence-preserves-continuity}, $f$ is continuous. For some $n$, $\norm{f-f_n}_{\infty}<1$, hence $f_n$ is bounded, thus $f\in C_b(X)$ and $\norm{f_n-f}_{\infty}\to0$.\\
\hlt{Structural Commentary}: Uniform Cauchy convergence builds the limit inside the same function space.\\
\hlt{Key Insight}: Completeness of $\C$ plus uniform convergence makes $C_b(X)$ complete.
\end{proof}
\end{theorem}

\begin{definition} \label{def:bounded-target-metric-function-space} \hlt{Bounded Target Metric for All Functions}\\
For metric spaces $(X,d_X)$ and $(Y,d_Y)$, let $\mathcal{F}=Y^X$ and set $\rho(y,z)\defeq\frac{d_Y(y,z)}{1+d_Y(y,z)}$ and $d_\rho(f,g)\defeq\sup_{x\in X}\rho(f(x),g(x))$.\\
Then $\rho<1$, so $d_\rho$ is finite on all of $\mathcal{F}$; $d_\rho$-convergence is uniform convergence measured by $d_Y$.
\end{definition}

\begin{theorem} \label{thm:bounded-target-metric-function-space-completeness} \hlt{Completeness and Closed Subspaces for $d_\rho$}\\
If the target metric space $(Y,d_Y)$ is complete, then:
\begin{enumerate}[label=(\roman*)]
\setlength{\itemsep}{0pt}
\item $(\mathcal{F},d_\rho)$ is complete;
\item $\mathcal{F}_b\defeq\{f\in\mathcal{F}:f(X)\text{ is }d_Y\text{-bounded}\}$ is closed in $\mathcal{F}$;
\item $C(X,Y)$ is closed in $\mathcal{F}$.
\end{enumerate}
\begin{proof}
\begin{enumerate}[label=(\roman*)]
\setlength{\itemsep}{0pt}
\item Let $(f_n)$ be $d_\rho$-Cauchy. For each $x\in X$, $(f_n(x))$ is Cauchy in $Y$; define $f(x)\defeq\lim_nf_n(x)$.\\
Given $\varepsilon>0$, choose $N$ with $d_\rho(f_n,f_m)<\varepsilon$ for $m,n\geq N$. Letting $m\to\infty$ gives $d_\rho(f_n,f)\leq\varepsilon$ for $n\geq N$.
\item If $f_n\in\mathcal{F}_b$ and $d_\rho(f_n,f)\to0$, choose $N$ with $d_\rho(f_N,f)<1/2$.\\
Then $d_Y(f_N(x),f(x))<1$ for all $x$. Since $f_N(X)$ is bounded, so is $f(X)$.
\item If $f_n\in C(X,Y)$ and $d_\rho(f_n,f)\to0$, then $f_n\to f$ uniformly in $d_Y$.\\
Given $x_0\in X$ and $\varepsilon>0$, choose $N$ with $d_Y(f_N(x),f(x))<\varepsilon/3$ for all $x$, then use continuity of $f_N$ at $x_0$ and the triangle inequality. Thus $f$ is continuous.
\end{enumerate}
\hlt{Structural Commentary}: Compress the target metric; then close the useful subspaces.\\
\hlt{Key Insight}: Completeness and uniform control make boundedness and continuity closed.
\end{proof}
\end{theorem}

\begin{definition} \label{def:pointwise-uniform-bounded-equicontinuous} \hlt{Boundedness and Equicontinuity of Function Families}\\
Let $\mathcal{F}$ be a family of functions from $E$ into a metric space $(Y,d_Y)$; in scalar cases, take $d_Y(u,v)=\abs{u-v}$.
\begin{enumerate}[label=(\roman*)]
\setlength{\itemsep}{0pt}
\item $(f_n)$ is \hlt{pointwise bounded} on $E$ if $(f_n(x))$ is bounded in $Y$ for each $x\in E$.
\item A scalar sequence $(f_n)$ is \hlt{uniformly bounded} on $E$ if $\exists M$ such that $\abs{f_n(x)}<M$ for all $x\in E,n\in\N$.
\item If $E$ lies in a metric space $(X,d_X)$, then $\mathcal{F}$ is \hlt{equicontinuous} on $E$ if $\forall\varepsilon>0$ $\exists\delta>0$,\\
such that $d_X(x,y)<\delta,\ x,y\in E,\ f\in\mathcal{F}$ $\Longrightarrow$ $d_Y(f(x),f(y))<\varepsilon$.
\item The family is \hlt{pointwise equicontinuous} if for every $p\in E$ and $\varepsilon>0$ there is $\delta=\delta(p,\varepsilon)>0$ such that $d_X(p,y)<\delta,\ y\in E,\ f\in\mathcal{F}$ implies $d_Y(f(p),f(y))<\varepsilon$.
\end{enumerate}
\end{definition}

\begin{figure}[H]
  \centering
  \begin{tikzpicture}[
  >=Stealth,
  x=0.82cm,
  y=0.86cm,
  every node/.style={font=\scriptsize},
  axis/.style={->, gray!70, line width=0.55pt},
  curvea/.style={blue!70!black, line width=0.9pt},
  curveb/.style={orange!85!black, line width=0.9pt},
  curvec/.style={teal!70!black, line width=0.9pt},
  guide/.style={gray!55, dashed, line width=0.5pt},
  domainband/.style={teal!10},
  epsbar/.style={black!65, line width=0.8pt},
  point/.style={circle, fill=black, inner sep=1.15pt}
]
  \draw[axis] (-0.35,0) -- (7.25,0) node[right] {$x$};
  \draw[axis] (0,-0.25) -- (0,3.7) node[above] {$y$};

  \path[domainband] (2.65,0) rectangle (4.35,3.35);
  \draw[guide] (2.65,0) -- (2.65,3.35);
  \draw[guide] (3.5,0) -- (3.5,3.35);
  \draw[guide] (4.35,0) -- (4.35,3.35);
  \node[below] at (2.65,0) {$p-\delta$};
  \node[below] at (3.5,0) {$p$};
  \node[below] at (4.35,0) {$p+\delta$};

  \draw[curvea] plot[smooth, tension=0.75] coordinates {
    (0.65,2.73) (1.55,2.88) (2.65,2.90) (3.5,2.82)
    (4.35,2.68) (5.45,2.56) (6.55,2.58)
  };
  \draw[curveb] plot[smooth, tension=0.75] coordinates {
    (0.65,1.86) (1.70,1.82) (2.65,1.83) (3.5,1.87)
    (4.35,1.95) (5.45,2.12) (6.55,2.20)
  };
  \draw[curvec] plot[smooth, tension=0.75] coordinates {
    (0.65,1.42) (1.65,1.43) (2.65,1.30) (3.5,1.11)
    (4.35,1.00) (5.45,0.85) (6.55,0.82)
  };

  \foreach \y in {2.82,1.87,1.11} {
    \node[point] at (3.5,\y) {};
  }

  \foreach \x/\y in {2.78/2.82,4.22/2.82,2.78/1.87,4.22/1.87,2.78/1.11,4.22/1.11} {
    \draw[epsbar] (\x,\y-0.18) -- (\x,\y+0.18);
  }

  \node[above] at (1.15,2.96) {$\mathcal{F}$};
  \node[black!65, right] at (4.36,2.82) {$\varepsilon$};
\end{tikzpicture}

  \caption{One input scale $\delta$ controls every curve in $\mathcal F$ within tolerance $\varepsilon$.}
  \label{fig:equicontinuity-family-schematic}
\end{figure}

\begin{theorem} \label{thm:pointwise-equicontinuity-compact-uniform} \hlt{Pointwise Equicontinuity Becomes Uniform on Compact Sets}\\
Every pointwise equicontinuous family $\mathcal{F}\subset C(K,\C)$ on compact $K$ is equicontinuous.
\begin{proof}
Otherwise, choose $\varepsilon>0$, $f_n\in\mathcal{F}$, and $x_n,t_n\in K$ with $d(x_n,t_n)<1/n$ but $\abs{f_n(x_n)-f_n(t_n)}\geq\varepsilon$.\\
By Theorem \ref{thm:compact-sequential-subsequence}, pass to $x_n\to x_0$; then $t_n\to x_0$ too.\\
Pointwise equicontinuity at $x_0$ gives $\abs{f_n(x_n)-f_n(t_n)}<\varepsilon$ for large $n$, a contradiction.\\
\hlt{Structural Commentary}: Compactness turns moving counterexamples into one base point.\\
\hlt{Key Insight}: The local $\delta(x,\varepsilon)$ becomes one uniform $\delta$ only after compactness prevents escape.
\end{proof}
\end{theorem}

\begin{theorem} \label{thm:uniform-convergence-implies-equicontinuity-on-compact} \hlt{Uniform Convergence Implies Equicontinuity on Compact Sets}\\
If $K$ is compact, $f_n\in C_b(K)$, and $f_n\rightrightarrows f$ on $K$, then $\{f_n:n\in\N\}$ is equicontinuous.
\begin{proof}
Choose $N$ so $\norm{f_n-f_N}_{\infty}<\varepsilon$ for $n>N$.\\
By Theorem \ref{thm:compact-continuity-uniform-continuity}, choose one $\delta>0$ giving $\varepsilon$-control for $f_1,\ldots,f_N$.\\
If $n>N$, $d(x,y)<\delta$, then $\abs{f_n(x)-f_n(y)} \leq \abs{f_n(x)-f_N(x)}+\abs{f_N(x)-f_N(y)}+\abs{f_N(y)-f_n(y)}<3\varepsilon$.\\
Together with the finite initial family, $\{f_n\}$ is equicontinuous.\\
\hlt{Structural Commentary}: A uniformly convergent tail stays close to one uniformly continuous function.\\
\hlt{Key Insight}: Compactness handles finitely many functions; uniform convergence handles the tail.
\end{proof}
\end{theorem}

\begin{example} \label{ex:uniform-boundedness-without-uniform-subsequence} \hlt{Uniform Boundedness Alone Is Not Enough}\\
On $[0,1]$, let $f_n(x)\defeq \frac{x^2}{x^2+(1-nx)^2}$. Then $\abs{f_n}\leq1$ and $f_n\to0$ pointwise, but $f_n(1/n)=1$; hence no uniformly convergent subsequence. The missing hypothesis is equicontinuity.
\end{example}

\begin{theorem} \label{thm:diagonal-selection-countable-pointwise} \hlt{Diagonal Selection on a Countable Set}\\
Every pointwise bounded sequence $(f_n)$ on countable $E$ has a subsequence converging pointwise on $E$.
\begin{proof}
List $E=\{x_1,x_2,\ldots\}$. Inductively choose nested subsequences $S_i$ so that $S_i(x_i)$ converges.\\
The diagonal subsequence is eventually contained in each $S_i$, hence converges at every $x_i$.\\
\hlt{Structural Commentary}: The diagonal keeps the $i$th convergence property after the $i$th step.\\
\hlt{Key Insight}: Countability permits diagonal convergence.
\end{proof}
\end{theorem}
